\documentclass[border=10pt,tikz]{standalone}
\usepackage{fontspec}
\setmainfont{Apple SD Gothic Neo}
\setsansfont{Apple SD Gothic Neo}
\usepackage{tikz}
\usetikzlibrary{positioning,arrows.meta,shapes,fit,backgrounds,calc}
\begin{document}
% _design.tex — shared TikZ design system for all blog diagrams
% Inputs nothing on its own — meant to be \input{} from each diagram .tex
% AFTER \begin{document}, BEFORE \begin{tikzpicture}.
%
% Companion files:
%   _design-math.tex — math diagrams (PGFPlots, book-notes palette)
%
% Provides a unified visual language across GoF / DSA / EMC++ / Beautiful / 시리즈:
%   - Colors: muted, accessible, color-coded by role
%   - Node styles: consistent sizing, rounded, subtle borders
%   - Edge styles: UML-style inheritance, aggregation, plus dashed alternates
%   - Annotation: callouts, notes, badges
%
% Usage:
%   \documentclass[border=10pt,tikz]{standalone}
%   \usepackage{tikz}
%   \usetikzlibrary{positioning,arrows.meta,shapes,fit,backgrounds,calc,decorations.pathreplacing}
%   \begin{document}
%   % _design.tex — shared TikZ design system for all blog diagrams
% Inputs nothing on its own — meant to be \input{} from each diagram .tex
% AFTER \begin{document}, BEFORE \begin{tikzpicture}.
%
% Companion files:
%   _design-math.tex — math diagrams (PGFPlots, book-notes palette)
%
% Provides a unified visual language across GoF / DSA / EMC++ / Beautiful / 시리즈:
%   - Colors: muted, accessible, color-coded by role
%   - Node styles: consistent sizing, rounded, subtle borders
%   - Edge styles: UML-style inheritance, aggregation, plus dashed alternates
%   - Annotation: callouts, notes, badges
%
% Usage:
%   \documentclass[border=10pt,tikz]{standalone}
%   \usepackage{tikz}
%   \usetikzlibrary{positioning,arrows.meta,shapes,fit,backgrounds,calc,decorations.pathreplacing}
%   \begin{document}
%   % _design.tex — shared TikZ design system for all blog diagrams
% Inputs nothing on its own — meant to be \input{} from each diagram .tex
% AFTER \begin{document}, BEFORE \begin{tikzpicture}.
%
% Companion files:
%   _design-math.tex — math diagrams (PGFPlots, book-notes palette)
%
% Provides a unified visual language across GoF / DSA / EMC++ / Beautiful / 시리즈:
%   - Colors: muted, accessible, color-coded by role
%   - Node styles: consistent sizing, rounded, subtle borders
%   - Edge styles: UML-style inheritance, aggregation, plus dashed alternates
%   - Annotation: callouts, notes, badges
%
% Usage:
%   \documentclass[border=10pt,tikz]{standalone}
%   \usepackage{tikz}
%   \usetikzlibrary{positioning,arrows.meta,shapes,fit,backgrounds,calc,decorations.pathreplacing}
%   \begin{document}
%   \input{../_design.tex}     % adjust relative path
%   \begin{tikzpicture}[blog]  % apply 'blog' style
%   ...
%
% Backward-compat note: This file only ADDS to the public API. Existing styles
% (abs/con/imp/cli/hi/vx/q/inh/agg/uses/arr/edge/alt/ret/thr/group/glabel) are
% preserved. New helpers are documented below.

% ─── Core palette (fills paired with darker borders) ─────────
\definecolor{absfill}{RGB}{220,232,247}   % cool blue
\definecolor{absbord}{RGB}{72,118,182}
\definecolor{confill}{RGB}{249,222,222}   % warm red
\definecolor{conbord}{RGB}{197,93,93}
\definecolor{impfill}{RGB}{223,238,222}   % soft green
\definecolor{impbord}{RGB}{99,154,93}
\definecolor{clifill}{RGB}{251,239,205}   % pale yellow
\definecolor{clibord}{RGB}{200,167,72}
\definecolor{hifill} {RGB}{253,224,200}   % accent orange
\definecolor{hibord} {RGB}{210,135,68}
\definecolor{nodefill}{RGB}{250,250,250}
\definecolor{nodebord}{RGB}{120,120,120}
\definecolor{edgec}  {RGB}{80,80,80}
\definecolor{edgesc} {RGB}{145,145,145}
\definecolor{groupbord}{RGB}{200,200,200}

% ─── Extended palette (new) ──────────────────────────────────
\definecolor{prpfill}{RGB}{233,224,247}   % muted purple (state / interface)
\definecolor{prpbord}{RGB}{132,103,182}
\definecolor{tealfill}{RGB}{215,235,235}  % teal (services / utility)
\definecolor{tealbord}{RGB}{ 78,135,140}
\definecolor{warnfill}{RGB}{251,229,229}  % deeper warn (error states)
\definecolor{warnbord}{RGB}{180, 55, 55}
\definecolor{okfill}  {RGB}{220,240,225}  % success
\definecolor{okbord}  {RGB}{ 75,140, 90}
\definecolor{neutfill}{RGB}{238,238,238}  % neutral panel
\definecolor{neutbord}{RGB}{170,170,170}
\definecolor{inkc}    {RGB}{ 30, 30, 30}  % strong ink (titles)
\definecolor{mutec}   {RGB}{120,120,120}  % muted text
\definecolor{focusc}  {RGB}{210, 70, 70}  % focus / highlight accent
\definecolor{gridc}   {RGB}{220,220,220}  % background grid

\tikzset{
  blog/.style={
    font=\sffamily\small\linespread{1.18}\selectfont,
    every node/.style={font=\sffamily\small\linespread{1.18}\selectfont}
  },
  % ─── existing node roles (unchanged) ───
  cls/.style={
    rectangle, rounded corners=2pt,
    draw=nodebord, line width=0.4pt,
    minimum height=8mm, minimum width=28mm,
    inner sep=3pt, fill=nodefill, align=center
  },
  abs/.style={cls, draw=absbord, fill=absfill, font=\sffamily\itshape\small},
  con/.style={cls, draw=conbord, fill=confill},
  imp/.style={cls, draw=impbord, fill=impfill},
  cli/.style={cls, draw=clibord, fill=clifill},
  hi/.style={cls, draw=hibord, fill=hifill, font=\sffamily\bfseries\small},
  % ─── new node roles ───
  prp/.style={cls, draw=prpbord, fill=prpfill},               % protocol / interface / state
  tealn/.style={cls, draw=tealbord, fill=tealfill},           % service / utility
  warn/.style={cls, draw=warnbord, fill=warnfill, font=\sffamily\bfseries\small},  % error state
  ok/.style={cls, draw=okbord, fill=okfill},                  % success state
  neut/.style={cls, draw=neutbord, fill=neutfill},            % neutral
  % ─── size variants ───
  clsSm/.style={cls, minimum height=6mm, minimum width=20mm, font=\sffamily\footnotesize, inner sep=2pt},
  clsLg/.style={cls, minimum height=11mm, minimum width=36mm, font=\sffamily\normalsize},
  absSm/.style={clsSm, draw=absbord, fill=absfill, font=\sffamily\itshape\footnotesize},
  conSm/.style={clsSm, draw=conbord, fill=confill},
  impSm/.style={clsSm, draw=impbord, fill=impfill},
  % ─── circle (existing) ───
  vx/.style={
    circle, draw=absbord, line width=0.4pt,
    minimum size=8mm, fill=absfill, inner sep=0pt
  },
  vxc/.style={vx, draw=conbord, fill=confill},
  vxh/.style={vx, draw=hibord, fill=hifill, font=\sffamily\bfseries},
  % new circle variants
  vxi/.style={vx, draw=impbord, fill=impfill},                % impl green
  vxp/.style={vx, draw=prpbord, fill=prpfill},                % protocol purple
  vxw/.style={vx, draw=warnbord, fill=warnfill},              % warn red
  vxn/.style={vx, draw=nodebord, fill=nodefill},              % neutral
  vxSm/.style={vx, minimum size=6mm, font=\sffamily\footnotesize},
  vxLg/.style={vx, minimum size=11mm, font=\sffamily\normalsize},
  % ─── decision diamond (existing) ───
  q/.style={
    diamond, draw=clibord, line width=0.4pt, fill=clifill,
    inner sep=2pt, aspect=2.4, align=center
  },
  % rounded callout / note (new)
  note/.style={
    rectangle, rounded corners=3pt,
    draw=clibord, line width=0.4pt, fill=clifill,
    inner sep=3pt, font=\sffamily\footnotesize, align=left
  },
  badge/.style={                                              % small inline marker
    rectangle, rounded corners=2pt,
    draw=hibord, line width=0.3pt, fill=hifill,
    inner sep=1.5pt, font=\sffamily\scriptsize\bfseries
  },
  % ─── edges (existing) ───
  inh/.style={                                  % UML inheritance
    -{Triangle[length=3mm,width=3mm,open]},
    draw=edgec, line width=0.5pt
  },
  agg/.style={                                  % UML aggregation
    {Diamond[length=2.5mm,width=2.5mm,open]}-{Stealth[length=2.2mm]},
    draw=edgec, line width=0.45pt
  },
  uses/.style={                                 % directed reference
    -{Stealth[length=2.2mm]}, draw=edgec, line width=0.5pt
  },
  arr/.style={                                  % plain arrow (graphs/trees)
    -{Stealth[length=2mm]}, draw=edgec, line width=0.45pt
  },
  edge/.style={                                 % undirected (trees)
    draw=edgec, line width=0.45pt
  },
  alt/.style={                                  % dashed: similar/alt
    -{Stealth[length=2mm]}, dashed, draw=edgesc, line width=0.45pt
  },
  ret/.style={                                  % dashed creates/returns
    -{Stealth[length=2mm]}, dashed, draw=edgesc, line width=0.45pt
  },
  thr/.style={                                  % threading link (red dashed)
    ->, dashed, draw=conbord, line width=0.5pt
  },
  % ─── new edges ───
  comp/.style={                                 % UML composition (filled diamond)
    {Diamond[length=2.5mm,width=2.5mm]}-{Stealth[length=2.2mm]},
    draw=edgec, line width=0.45pt
  },
  realize/.style={                              % UML realize (impl interface)
    -{Triangle[length=3mm,width=3mm,open]}, dashed,
    draw=edgec, line width=0.5pt
  },
  msgsync/.style={                              % synchronous message (sequence)
    -{Stealth[length=2.6mm]}, draw=edgec, line width=0.55pt
  },
  msgasync/.style={                             % asynchronous message
    -{Stealth[length=2.4mm,inset=0.4mm,fill=edgec]},
    draw=edgec, line width=0.5pt
  },
  msgret/.style={                               % return message (dashed)
    -{Stealth[length=2.2mm]}, dashed, draw=edgesc, line width=0.45pt
  },
  flow/.style={                                 % control flow (thicker, hi-vis)
    -{Stealth[length=2.4mm]}, draw=inkc, line width=0.7pt
  },
  focus/.style={                                % focus / hot path
    -{Stealth[length=2.4mm]}, draw=focusc, line width=0.8pt
  },
  selfloop/.style={                             % self loop on a node
    ->, draw=edgec, line width=0.45pt, in=120, out=60, loop, looseness=8
  },
  % ─── group / region (existing + new) ───
  group/.style={
    rectangle, rounded corners=4pt,
    draw=groupbord, line width=0.4pt,
    inner sep=10pt, fill=black!2
  },
  glabel/.style={font=\sffamily\bfseries\small, anchor=north west, inner sep=2pt},
  panel/.style={                                % stronger panel for layered diagrams
    rectangle, rounded corners=4pt,
    draw=neutbord, line width=0.45pt,
    inner sep=10pt, fill=neutfill!50
  },
  panelTitle/.style={font=\sffamily\bfseries\small, text=inkc, anchor=north west, inner sep=2pt},
  panelSub/.style={font=\sffamily\footnotesize, text=mutec, anchor=north west, inner sep=2pt},
  % ─── typography ───
  caption/.style={font=\sffamily\footnotesize\itshape, text=mutec, align=center},
  legendkey/.style={rectangle, minimum width=10mm, minimum height=3mm, inner sep=0pt, line width=0.3pt},
  % ─── grid / coordinate hints ───
  bggrid/.style={draw=gridc, line width=0.3pt, dotted},
}

% Korean font convenience macro — included by every Hangul diagram
\providecommand{\KOR}{}
     % adjust relative path
%   \begin{tikzpicture}[blog]  % apply 'blog' style
%   ...
%
% Backward-compat note: This file only ADDS to the public API. Existing styles
% (abs/con/imp/cli/hi/vx/q/inh/agg/uses/arr/edge/alt/ret/thr/group/glabel) are
% preserved. New helpers are documented below.

% ─── Core palette (fills paired with darker borders) ─────────
\definecolor{absfill}{RGB}{220,232,247}   % cool blue
\definecolor{absbord}{RGB}{72,118,182}
\definecolor{confill}{RGB}{249,222,222}   % warm red
\definecolor{conbord}{RGB}{197,93,93}
\definecolor{impfill}{RGB}{223,238,222}   % soft green
\definecolor{impbord}{RGB}{99,154,93}
\definecolor{clifill}{RGB}{251,239,205}   % pale yellow
\definecolor{clibord}{RGB}{200,167,72}
\definecolor{hifill} {RGB}{253,224,200}   % accent orange
\definecolor{hibord} {RGB}{210,135,68}
\definecolor{nodefill}{RGB}{250,250,250}
\definecolor{nodebord}{RGB}{120,120,120}
\definecolor{edgec}  {RGB}{80,80,80}
\definecolor{edgesc} {RGB}{145,145,145}
\definecolor{groupbord}{RGB}{200,200,200}

% ─── Extended palette (new) ──────────────────────────────────
\definecolor{prpfill}{RGB}{233,224,247}   % muted purple (state / interface)
\definecolor{prpbord}{RGB}{132,103,182}
\definecolor{tealfill}{RGB}{215,235,235}  % teal (services / utility)
\definecolor{tealbord}{RGB}{ 78,135,140}
\definecolor{warnfill}{RGB}{251,229,229}  % deeper warn (error states)
\definecolor{warnbord}{RGB}{180, 55, 55}
\definecolor{okfill}  {RGB}{220,240,225}  % success
\definecolor{okbord}  {RGB}{ 75,140, 90}
\definecolor{neutfill}{RGB}{238,238,238}  % neutral panel
\definecolor{neutbord}{RGB}{170,170,170}
\definecolor{inkc}    {RGB}{ 30, 30, 30}  % strong ink (titles)
\definecolor{mutec}   {RGB}{120,120,120}  % muted text
\definecolor{focusc}  {RGB}{210, 70, 70}  % focus / highlight accent
\definecolor{gridc}   {RGB}{220,220,220}  % background grid

\tikzset{
  blog/.style={
    font=\sffamily\small\linespread{1.18}\selectfont,
    every node/.style={font=\sffamily\small\linespread{1.18}\selectfont}
  },
  % ─── existing node roles (unchanged) ───
  cls/.style={
    rectangle, rounded corners=2pt,
    draw=nodebord, line width=0.4pt,
    minimum height=8mm, minimum width=28mm,
    inner sep=3pt, fill=nodefill, align=center
  },
  abs/.style={cls, draw=absbord, fill=absfill, font=\sffamily\itshape\small},
  con/.style={cls, draw=conbord, fill=confill},
  imp/.style={cls, draw=impbord, fill=impfill},
  cli/.style={cls, draw=clibord, fill=clifill},
  hi/.style={cls, draw=hibord, fill=hifill, font=\sffamily\bfseries\small},
  % ─── new node roles ───
  prp/.style={cls, draw=prpbord, fill=prpfill},               % protocol / interface / state
  tealn/.style={cls, draw=tealbord, fill=tealfill},           % service / utility
  warn/.style={cls, draw=warnbord, fill=warnfill, font=\sffamily\bfseries\small},  % error state
  ok/.style={cls, draw=okbord, fill=okfill},                  % success state
  neut/.style={cls, draw=neutbord, fill=neutfill},            % neutral
  % ─── size variants ───
  clsSm/.style={cls, minimum height=6mm, minimum width=20mm, font=\sffamily\footnotesize, inner sep=2pt},
  clsLg/.style={cls, minimum height=11mm, minimum width=36mm, font=\sffamily\normalsize},
  absSm/.style={clsSm, draw=absbord, fill=absfill, font=\sffamily\itshape\footnotesize},
  conSm/.style={clsSm, draw=conbord, fill=confill},
  impSm/.style={clsSm, draw=impbord, fill=impfill},
  % ─── circle (existing) ───
  vx/.style={
    circle, draw=absbord, line width=0.4pt,
    minimum size=8mm, fill=absfill, inner sep=0pt
  },
  vxc/.style={vx, draw=conbord, fill=confill},
  vxh/.style={vx, draw=hibord, fill=hifill, font=\sffamily\bfseries},
  % new circle variants
  vxi/.style={vx, draw=impbord, fill=impfill},                % impl green
  vxp/.style={vx, draw=prpbord, fill=prpfill},                % protocol purple
  vxw/.style={vx, draw=warnbord, fill=warnfill},              % warn red
  vxn/.style={vx, draw=nodebord, fill=nodefill},              % neutral
  vxSm/.style={vx, minimum size=6mm, font=\sffamily\footnotesize},
  vxLg/.style={vx, minimum size=11mm, font=\sffamily\normalsize},
  % ─── decision diamond (existing) ───
  q/.style={
    diamond, draw=clibord, line width=0.4pt, fill=clifill,
    inner sep=2pt, aspect=2.4, align=center
  },
  % rounded callout / note (new)
  note/.style={
    rectangle, rounded corners=3pt,
    draw=clibord, line width=0.4pt, fill=clifill,
    inner sep=3pt, font=\sffamily\footnotesize, align=left
  },
  badge/.style={                                              % small inline marker
    rectangle, rounded corners=2pt,
    draw=hibord, line width=0.3pt, fill=hifill,
    inner sep=1.5pt, font=\sffamily\scriptsize\bfseries
  },
  % ─── edges (existing) ───
  inh/.style={                                  % UML inheritance
    -{Triangle[length=3mm,width=3mm,open]},
    draw=edgec, line width=0.5pt
  },
  agg/.style={                                  % UML aggregation
    {Diamond[length=2.5mm,width=2.5mm,open]}-{Stealth[length=2.2mm]},
    draw=edgec, line width=0.45pt
  },
  uses/.style={                                 % directed reference
    -{Stealth[length=2.2mm]}, draw=edgec, line width=0.5pt
  },
  arr/.style={                                  % plain arrow (graphs/trees)
    -{Stealth[length=2mm]}, draw=edgec, line width=0.45pt
  },
  edge/.style={                                 % undirected (trees)
    draw=edgec, line width=0.45pt
  },
  alt/.style={                                  % dashed: similar/alt
    -{Stealth[length=2mm]}, dashed, draw=edgesc, line width=0.45pt
  },
  ret/.style={                                  % dashed creates/returns
    -{Stealth[length=2mm]}, dashed, draw=edgesc, line width=0.45pt
  },
  thr/.style={                                  % threading link (red dashed)
    ->, dashed, draw=conbord, line width=0.5pt
  },
  % ─── new edges ───
  comp/.style={                                 % UML composition (filled diamond)
    {Diamond[length=2.5mm,width=2.5mm]}-{Stealth[length=2.2mm]},
    draw=edgec, line width=0.45pt
  },
  realize/.style={                              % UML realize (impl interface)
    -{Triangle[length=3mm,width=3mm,open]}, dashed,
    draw=edgec, line width=0.5pt
  },
  msgsync/.style={                              % synchronous message (sequence)
    -{Stealth[length=2.6mm]}, draw=edgec, line width=0.55pt
  },
  msgasync/.style={                             % asynchronous message
    -{Stealth[length=2.4mm,inset=0.4mm,fill=edgec]},
    draw=edgec, line width=0.5pt
  },
  msgret/.style={                               % return message (dashed)
    -{Stealth[length=2.2mm]}, dashed, draw=edgesc, line width=0.45pt
  },
  flow/.style={                                 % control flow (thicker, hi-vis)
    -{Stealth[length=2.4mm]}, draw=inkc, line width=0.7pt
  },
  focus/.style={                                % focus / hot path
    -{Stealth[length=2.4mm]}, draw=focusc, line width=0.8pt
  },
  selfloop/.style={                             % self loop on a node
    ->, draw=edgec, line width=0.45pt, in=120, out=60, loop, looseness=8
  },
  % ─── group / region (existing + new) ───
  group/.style={
    rectangle, rounded corners=4pt,
    draw=groupbord, line width=0.4pt,
    inner sep=10pt, fill=black!2
  },
  glabel/.style={font=\sffamily\bfseries\small, anchor=north west, inner sep=2pt},
  panel/.style={                                % stronger panel for layered diagrams
    rectangle, rounded corners=4pt,
    draw=neutbord, line width=0.45pt,
    inner sep=10pt, fill=neutfill!50
  },
  panelTitle/.style={font=\sffamily\bfseries\small, text=inkc, anchor=north west, inner sep=2pt},
  panelSub/.style={font=\sffamily\footnotesize, text=mutec, anchor=north west, inner sep=2pt},
  % ─── typography ───
  caption/.style={font=\sffamily\footnotesize\itshape, text=mutec, align=center},
  legendkey/.style={rectangle, minimum width=10mm, minimum height=3mm, inner sep=0pt, line width=0.3pt},
  % ─── grid / coordinate hints ───
  bggrid/.style={draw=gridc, line width=0.3pt, dotted},
}

% Korean font convenience macro — included by every Hangul diagram
\providecommand{\KOR}{}
     % adjust relative path
%   \begin{tikzpicture}[blog]  % apply 'blog' style
%   ...
%
% Backward-compat note: This file only ADDS to the public API. Existing styles
% (abs/con/imp/cli/hi/vx/q/inh/agg/uses/arr/edge/alt/ret/thr/group/glabel) are
% preserved. New helpers are documented below.

% ─── Core palette (fills paired with darker borders) ─────────
\definecolor{absfill}{RGB}{220,232,247}   % cool blue
\definecolor{absbord}{RGB}{72,118,182}
\definecolor{confill}{RGB}{249,222,222}   % warm red
\definecolor{conbord}{RGB}{197,93,93}
\definecolor{impfill}{RGB}{223,238,222}   % soft green
\definecolor{impbord}{RGB}{99,154,93}
\definecolor{clifill}{RGB}{251,239,205}   % pale yellow
\definecolor{clibord}{RGB}{200,167,72}
\definecolor{hifill} {RGB}{253,224,200}   % accent orange
\definecolor{hibord} {RGB}{210,135,68}
\definecolor{nodefill}{RGB}{250,250,250}
\definecolor{nodebord}{RGB}{120,120,120}
\definecolor{edgec}  {RGB}{80,80,80}
\definecolor{edgesc} {RGB}{145,145,145}
\definecolor{groupbord}{RGB}{200,200,200}

% ─── Extended palette (new) ──────────────────────────────────
\definecolor{prpfill}{RGB}{233,224,247}   % muted purple (state / interface)
\definecolor{prpbord}{RGB}{132,103,182}
\definecolor{tealfill}{RGB}{215,235,235}  % teal (services / utility)
\definecolor{tealbord}{RGB}{ 78,135,140}
\definecolor{warnfill}{RGB}{251,229,229}  % deeper warn (error states)
\definecolor{warnbord}{RGB}{180, 55, 55}
\definecolor{okfill}  {RGB}{220,240,225}  % success
\definecolor{okbord}  {RGB}{ 75,140, 90}
\definecolor{neutfill}{RGB}{238,238,238}  % neutral panel
\definecolor{neutbord}{RGB}{170,170,170}
\definecolor{inkc}    {RGB}{ 30, 30, 30}  % strong ink (titles)
\definecolor{mutec}   {RGB}{120,120,120}  % muted text
\definecolor{focusc}  {RGB}{210, 70, 70}  % focus / highlight accent
\definecolor{gridc}   {RGB}{220,220,220}  % background grid

\tikzset{
  blog/.style={
    font=\sffamily\small\linespread{1.18}\selectfont,
    every node/.style={font=\sffamily\small\linespread{1.18}\selectfont}
  },
  % ─── existing node roles (unchanged) ───
  cls/.style={
    rectangle, rounded corners=2pt,
    draw=nodebord, line width=0.4pt,
    minimum height=8mm, minimum width=28mm,
    inner sep=3pt, fill=nodefill, align=center
  },
  abs/.style={cls, draw=absbord, fill=absfill, font=\sffamily\itshape\small},
  con/.style={cls, draw=conbord, fill=confill},
  imp/.style={cls, draw=impbord, fill=impfill},
  cli/.style={cls, draw=clibord, fill=clifill},
  hi/.style={cls, draw=hibord, fill=hifill, font=\sffamily\bfseries\small},
  % ─── new node roles ───
  prp/.style={cls, draw=prpbord, fill=prpfill},               % protocol / interface / state
  tealn/.style={cls, draw=tealbord, fill=tealfill},           % service / utility
  warn/.style={cls, draw=warnbord, fill=warnfill, font=\sffamily\bfseries\small},  % error state
  ok/.style={cls, draw=okbord, fill=okfill},                  % success state
  neut/.style={cls, draw=neutbord, fill=neutfill},            % neutral
  % ─── size variants ───
  clsSm/.style={cls, minimum height=6mm, minimum width=20mm, font=\sffamily\footnotesize, inner sep=2pt},
  clsLg/.style={cls, minimum height=11mm, minimum width=36mm, font=\sffamily\normalsize},
  absSm/.style={clsSm, draw=absbord, fill=absfill, font=\sffamily\itshape\footnotesize},
  conSm/.style={clsSm, draw=conbord, fill=confill},
  impSm/.style={clsSm, draw=impbord, fill=impfill},
  % ─── circle (existing) ───
  vx/.style={
    circle, draw=absbord, line width=0.4pt,
    minimum size=8mm, fill=absfill, inner sep=0pt
  },
  vxc/.style={vx, draw=conbord, fill=confill},
  vxh/.style={vx, draw=hibord, fill=hifill, font=\sffamily\bfseries},
  % new circle variants
  vxi/.style={vx, draw=impbord, fill=impfill},                % impl green
  vxp/.style={vx, draw=prpbord, fill=prpfill},                % protocol purple
  vxw/.style={vx, draw=warnbord, fill=warnfill},              % warn red
  vxn/.style={vx, draw=nodebord, fill=nodefill},              % neutral
  vxSm/.style={vx, minimum size=6mm, font=\sffamily\footnotesize},
  vxLg/.style={vx, minimum size=11mm, font=\sffamily\normalsize},
  % ─── decision diamond (existing) ───
  q/.style={
    diamond, draw=clibord, line width=0.4pt, fill=clifill,
    inner sep=2pt, aspect=2.4, align=center
  },
  % rounded callout / note (new)
  note/.style={
    rectangle, rounded corners=3pt,
    draw=clibord, line width=0.4pt, fill=clifill,
    inner sep=3pt, font=\sffamily\footnotesize, align=left
  },
  badge/.style={                                              % small inline marker
    rectangle, rounded corners=2pt,
    draw=hibord, line width=0.3pt, fill=hifill,
    inner sep=1.5pt, font=\sffamily\scriptsize\bfseries
  },
  % ─── edges (existing) ───
  inh/.style={                                  % UML inheritance
    -{Triangle[length=3mm,width=3mm,open]},
    draw=edgec, line width=0.5pt
  },
  agg/.style={                                  % UML aggregation
    {Diamond[length=2.5mm,width=2.5mm,open]}-{Stealth[length=2.2mm]},
    draw=edgec, line width=0.45pt
  },
  uses/.style={                                 % directed reference
    -{Stealth[length=2.2mm]}, draw=edgec, line width=0.5pt
  },
  arr/.style={                                  % plain arrow (graphs/trees)
    -{Stealth[length=2mm]}, draw=edgec, line width=0.45pt
  },
  edge/.style={                                 % undirected (trees)
    draw=edgec, line width=0.45pt
  },
  alt/.style={                                  % dashed: similar/alt
    -{Stealth[length=2mm]}, dashed, draw=edgesc, line width=0.45pt
  },
  ret/.style={                                  % dashed creates/returns
    -{Stealth[length=2mm]}, dashed, draw=edgesc, line width=0.45pt
  },
  thr/.style={                                  % threading link (red dashed)
    ->, dashed, draw=conbord, line width=0.5pt
  },
  % ─── new edges ───
  comp/.style={                                 % UML composition (filled diamond)
    {Diamond[length=2.5mm,width=2.5mm]}-{Stealth[length=2.2mm]},
    draw=edgec, line width=0.45pt
  },
  realize/.style={                              % UML realize (impl interface)
    -{Triangle[length=3mm,width=3mm,open]}, dashed,
    draw=edgec, line width=0.5pt
  },
  msgsync/.style={                              % synchronous message (sequence)
    -{Stealth[length=2.6mm]}, draw=edgec, line width=0.55pt
  },
  msgasync/.style={                             % asynchronous message
    -{Stealth[length=2.4mm,inset=0.4mm,fill=edgec]},
    draw=edgec, line width=0.5pt
  },
  msgret/.style={                               % return message (dashed)
    -{Stealth[length=2.2mm]}, dashed, draw=edgesc, line width=0.45pt
  },
  flow/.style={                                 % control flow (thicker, hi-vis)
    -{Stealth[length=2.4mm]}, draw=inkc, line width=0.7pt
  },
  focus/.style={                                % focus / hot path
    -{Stealth[length=2.4mm]}, draw=focusc, line width=0.8pt
  },
  selfloop/.style={                             % self loop on a node
    ->, draw=edgec, line width=0.45pt, in=120, out=60, loop, looseness=8
  },
  % ─── group / region (existing + new) ───
  group/.style={
    rectangle, rounded corners=4pt,
    draw=groupbord, line width=0.4pt,
    inner sep=10pt, fill=black!2
  },
  glabel/.style={font=\sffamily\bfseries\small, anchor=north west, inner sep=2pt},
  panel/.style={                                % stronger panel for layered diagrams
    rectangle, rounded corners=4pt,
    draw=neutbord, line width=0.45pt,
    inner sep=10pt, fill=neutfill!50
  },
  panelTitle/.style={font=\sffamily\bfseries\small, text=inkc, anchor=north west, inner sep=2pt},
  panelSub/.style={font=\sffamily\footnotesize, text=mutec, anchor=north west, inner sep=2pt},
  % ─── typography ───
  caption/.style={font=\sffamily\footnotesize\itshape, text=mutec, align=center},
  legendkey/.style={rectangle, minimum width=10mm, minimum height=3mm, inner sep=0pt, line width=0.3pt},
  % ─── grid / coordinate hints ───
  bggrid/.style={draw=gridc, line width=0.3pt, dotted},
}

% Korean font convenience macro — included by every Hangul diagram
\providecommand{\KOR}{}

% _design-sequence.tex — TikZ sequence diagram styles
%
% Mermaid는 *간단한* sequence diagram에 좋지만, 복잡한 경우 (여러 actor,
% 중첩 activation, alt/opt/loop frame, 자세한 라벨)는 auto-layout이 들쭉날쭉.
% 정밀 좌표가 필요할 때 TikZ로 그리는 표준.
%
% Companion to _design.tex — same palette and ink colors.
%
% Usage:
%   \documentclass[border=10pt,tikz]{standalone}
%   \usepackage{tikz}
%   \usetikzlibrary{positioning,arrows.meta,shapes,fit,calc,decorations.pathreplacing}
%   \begin{document}
%   % _design.tex — shared TikZ design system for all blog diagrams
% Inputs nothing on its own — meant to be \input{} from each diagram .tex
% AFTER \begin{document}, BEFORE \begin{tikzpicture}.
%
% Companion files:
%   _design-math.tex — math diagrams (PGFPlots, book-notes palette)
%
% Provides a unified visual language across GoF / DSA / EMC++ / Beautiful / 시리즈:
%   - Colors: muted, accessible, color-coded by role
%   - Node styles: consistent sizing, rounded, subtle borders
%   - Edge styles: UML-style inheritance, aggregation, plus dashed alternates
%   - Annotation: callouts, notes, badges
%
% Usage:
%   \documentclass[border=10pt,tikz]{standalone}
%   \usepackage{tikz}
%   \usetikzlibrary{positioning,arrows.meta,shapes,fit,backgrounds,calc,decorations.pathreplacing}
%   \begin{document}
%   % _design.tex — shared TikZ design system for all blog diagrams
% Inputs nothing on its own — meant to be \input{} from each diagram .tex
% AFTER \begin{document}, BEFORE \begin{tikzpicture}.
%
% Companion files:
%   _design-math.tex — math diagrams (PGFPlots, book-notes palette)
%
% Provides a unified visual language across GoF / DSA / EMC++ / Beautiful / 시리즈:
%   - Colors: muted, accessible, color-coded by role
%   - Node styles: consistent sizing, rounded, subtle borders
%   - Edge styles: UML-style inheritance, aggregation, plus dashed alternates
%   - Annotation: callouts, notes, badges
%
% Usage:
%   \documentclass[border=10pt,tikz]{standalone}
%   \usepackage{tikz}
%   \usetikzlibrary{positioning,arrows.meta,shapes,fit,backgrounds,calc,decorations.pathreplacing}
%   \begin{document}
%   \input{../_design.tex}     % adjust relative path
%   \begin{tikzpicture}[blog]  % apply 'blog' style
%   ...
%
% Backward-compat note: This file only ADDS to the public API. Existing styles
% (abs/con/imp/cli/hi/vx/q/inh/agg/uses/arr/edge/alt/ret/thr/group/glabel) are
% preserved. New helpers are documented below.

% ─── Core palette (fills paired with darker borders) ─────────
\definecolor{absfill}{RGB}{220,232,247}   % cool blue
\definecolor{absbord}{RGB}{72,118,182}
\definecolor{confill}{RGB}{249,222,222}   % warm red
\definecolor{conbord}{RGB}{197,93,93}
\definecolor{impfill}{RGB}{223,238,222}   % soft green
\definecolor{impbord}{RGB}{99,154,93}
\definecolor{clifill}{RGB}{251,239,205}   % pale yellow
\definecolor{clibord}{RGB}{200,167,72}
\definecolor{hifill} {RGB}{253,224,200}   % accent orange
\definecolor{hibord} {RGB}{210,135,68}
\definecolor{nodefill}{RGB}{250,250,250}
\definecolor{nodebord}{RGB}{120,120,120}
\definecolor{edgec}  {RGB}{80,80,80}
\definecolor{edgesc} {RGB}{145,145,145}
\definecolor{groupbord}{RGB}{200,200,200}

% ─── Extended palette (new) ──────────────────────────────────
\definecolor{prpfill}{RGB}{233,224,247}   % muted purple (state / interface)
\definecolor{prpbord}{RGB}{132,103,182}
\definecolor{tealfill}{RGB}{215,235,235}  % teal (services / utility)
\definecolor{tealbord}{RGB}{ 78,135,140}
\definecolor{warnfill}{RGB}{251,229,229}  % deeper warn (error states)
\definecolor{warnbord}{RGB}{180, 55, 55}
\definecolor{okfill}  {RGB}{220,240,225}  % success
\definecolor{okbord}  {RGB}{ 75,140, 90}
\definecolor{neutfill}{RGB}{238,238,238}  % neutral panel
\definecolor{neutbord}{RGB}{170,170,170}
\definecolor{inkc}    {RGB}{ 30, 30, 30}  % strong ink (titles)
\definecolor{mutec}   {RGB}{120,120,120}  % muted text
\definecolor{focusc}  {RGB}{210, 70, 70}  % focus / highlight accent
\definecolor{gridc}   {RGB}{220,220,220}  % background grid

\tikzset{
  blog/.style={
    font=\sffamily\small\linespread{1.18}\selectfont,
    every node/.style={font=\sffamily\small\linespread{1.18}\selectfont}
  },
  % ─── existing node roles (unchanged) ───
  cls/.style={
    rectangle, rounded corners=2pt,
    draw=nodebord, line width=0.4pt,
    minimum height=8mm, minimum width=28mm,
    inner sep=3pt, fill=nodefill, align=center
  },
  abs/.style={cls, draw=absbord, fill=absfill, font=\sffamily\itshape\small},
  con/.style={cls, draw=conbord, fill=confill},
  imp/.style={cls, draw=impbord, fill=impfill},
  cli/.style={cls, draw=clibord, fill=clifill},
  hi/.style={cls, draw=hibord, fill=hifill, font=\sffamily\bfseries\small},
  % ─── new node roles ───
  prp/.style={cls, draw=prpbord, fill=prpfill},               % protocol / interface / state
  tealn/.style={cls, draw=tealbord, fill=tealfill},           % service / utility
  warn/.style={cls, draw=warnbord, fill=warnfill, font=\sffamily\bfseries\small},  % error state
  ok/.style={cls, draw=okbord, fill=okfill},                  % success state
  neut/.style={cls, draw=neutbord, fill=neutfill},            % neutral
  % ─── size variants ───
  clsSm/.style={cls, minimum height=6mm, minimum width=20mm, font=\sffamily\footnotesize, inner sep=2pt},
  clsLg/.style={cls, minimum height=11mm, minimum width=36mm, font=\sffamily\normalsize},
  absSm/.style={clsSm, draw=absbord, fill=absfill, font=\sffamily\itshape\footnotesize},
  conSm/.style={clsSm, draw=conbord, fill=confill},
  impSm/.style={clsSm, draw=impbord, fill=impfill},
  % ─── circle (existing) ───
  vx/.style={
    circle, draw=absbord, line width=0.4pt,
    minimum size=8mm, fill=absfill, inner sep=0pt
  },
  vxc/.style={vx, draw=conbord, fill=confill},
  vxh/.style={vx, draw=hibord, fill=hifill, font=\sffamily\bfseries},
  % new circle variants
  vxi/.style={vx, draw=impbord, fill=impfill},                % impl green
  vxp/.style={vx, draw=prpbord, fill=prpfill},                % protocol purple
  vxw/.style={vx, draw=warnbord, fill=warnfill},              % warn red
  vxn/.style={vx, draw=nodebord, fill=nodefill},              % neutral
  vxSm/.style={vx, minimum size=6mm, font=\sffamily\footnotesize},
  vxLg/.style={vx, minimum size=11mm, font=\sffamily\normalsize},
  % ─── decision diamond (existing) ───
  q/.style={
    diamond, draw=clibord, line width=0.4pt, fill=clifill,
    inner sep=2pt, aspect=2.4, align=center
  },
  % rounded callout / note (new)
  note/.style={
    rectangle, rounded corners=3pt,
    draw=clibord, line width=0.4pt, fill=clifill,
    inner sep=3pt, font=\sffamily\footnotesize, align=left
  },
  badge/.style={                                              % small inline marker
    rectangle, rounded corners=2pt,
    draw=hibord, line width=0.3pt, fill=hifill,
    inner sep=1.5pt, font=\sffamily\scriptsize\bfseries
  },
  % ─── edges (existing) ───
  inh/.style={                                  % UML inheritance
    -{Triangle[length=3mm,width=3mm,open]},
    draw=edgec, line width=0.5pt
  },
  agg/.style={                                  % UML aggregation
    {Diamond[length=2.5mm,width=2.5mm,open]}-{Stealth[length=2.2mm]},
    draw=edgec, line width=0.45pt
  },
  uses/.style={                                 % directed reference
    -{Stealth[length=2.2mm]}, draw=edgec, line width=0.5pt
  },
  arr/.style={                                  % plain arrow (graphs/trees)
    -{Stealth[length=2mm]}, draw=edgec, line width=0.45pt
  },
  edge/.style={                                 % undirected (trees)
    draw=edgec, line width=0.45pt
  },
  alt/.style={                                  % dashed: similar/alt
    -{Stealth[length=2mm]}, dashed, draw=edgesc, line width=0.45pt
  },
  ret/.style={                                  % dashed creates/returns
    -{Stealth[length=2mm]}, dashed, draw=edgesc, line width=0.45pt
  },
  thr/.style={                                  % threading link (red dashed)
    ->, dashed, draw=conbord, line width=0.5pt
  },
  % ─── new edges ───
  comp/.style={                                 % UML composition (filled diamond)
    {Diamond[length=2.5mm,width=2.5mm]}-{Stealth[length=2.2mm]},
    draw=edgec, line width=0.45pt
  },
  realize/.style={                              % UML realize (impl interface)
    -{Triangle[length=3mm,width=3mm,open]}, dashed,
    draw=edgec, line width=0.5pt
  },
  msgsync/.style={                              % synchronous message (sequence)
    -{Stealth[length=2.6mm]}, draw=edgec, line width=0.55pt
  },
  msgasync/.style={                             % asynchronous message
    -{Stealth[length=2.4mm,inset=0.4mm,fill=edgec]},
    draw=edgec, line width=0.5pt
  },
  msgret/.style={                               % return message (dashed)
    -{Stealth[length=2.2mm]}, dashed, draw=edgesc, line width=0.45pt
  },
  flow/.style={                                 % control flow (thicker, hi-vis)
    -{Stealth[length=2.4mm]}, draw=inkc, line width=0.7pt
  },
  focus/.style={                                % focus / hot path
    -{Stealth[length=2.4mm]}, draw=focusc, line width=0.8pt
  },
  selfloop/.style={                             % self loop on a node
    ->, draw=edgec, line width=0.45pt, in=120, out=60, loop, looseness=8
  },
  % ─── group / region (existing + new) ───
  group/.style={
    rectangle, rounded corners=4pt,
    draw=groupbord, line width=0.4pt,
    inner sep=10pt, fill=black!2
  },
  glabel/.style={font=\sffamily\bfseries\small, anchor=north west, inner sep=2pt},
  panel/.style={                                % stronger panel for layered diagrams
    rectangle, rounded corners=4pt,
    draw=neutbord, line width=0.45pt,
    inner sep=10pt, fill=neutfill!50
  },
  panelTitle/.style={font=\sffamily\bfseries\small, text=inkc, anchor=north west, inner sep=2pt},
  panelSub/.style={font=\sffamily\footnotesize, text=mutec, anchor=north west, inner sep=2pt},
  % ─── typography ───
  caption/.style={font=\sffamily\footnotesize\itshape, text=mutec, align=center},
  legendkey/.style={rectangle, minimum width=10mm, minimum height=3mm, inner sep=0pt, line width=0.3pt},
  % ─── grid / coordinate hints ───
  bggrid/.style={draw=gridc, line width=0.3pt, dotted},
}

% Korean font convenience macro — included by every Hangul diagram
\providecommand{\KOR}{}
     % adjust relative path
%   \begin{tikzpicture}[blog]  % apply 'blog' style
%   ...
%
% Backward-compat note: This file only ADDS to the public API. Existing styles
% (abs/con/imp/cli/hi/vx/q/inh/agg/uses/arr/edge/alt/ret/thr/group/glabel) are
% preserved. New helpers are documented below.

% ─── Core palette (fills paired with darker borders) ─────────
\definecolor{absfill}{RGB}{220,232,247}   % cool blue
\definecolor{absbord}{RGB}{72,118,182}
\definecolor{confill}{RGB}{249,222,222}   % warm red
\definecolor{conbord}{RGB}{197,93,93}
\definecolor{impfill}{RGB}{223,238,222}   % soft green
\definecolor{impbord}{RGB}{99,154,93}
\definecolor{clifill}{RGB}{251,239,205}   % pale yellow
\definecolor{clibord}{RGB}{200,167,72}
\definecolor{hifill} {RGB}{253,224,200}   % accent orange
\definecolor{hibord} {RGB}{210,135,68}
\definecolor{nodefill}{RGB}{250,250,250}
\definecolor{nodebord}{RGB}{120,120,120}
\definecolor{edgec}  {RGB}{80,80,80}
\definecolor{edgesc} {RGB}{145,145,145}
\definecolor{groupbord}{RGB}{200,200,200}

% ─── Extended palette (new) ──────────────────────────────────
\definecolor{prpfill}{RGB}{233,224,247}   % muted purple (state / interface)
\definecolor{prpbord}{RGB}{132,103,182}
\definecolor{tealfill}{RGB}{215,235,235}  % teal (services / utility)
\definecolor{tealbord}{RGB}{ 78,135,140}
\definecolor{warnfill}{RGB}{251,229,229}  % deeper warn (error states)
\definecolor{warnbord}{RGB}{180, 55, 55}
\definecolor{okfill}  {RGB}{220,240,225}  % success
\definecolor{okbord}  {RGB}{ 75,140, 90}
\definecolor{neutfill}{RGB}{238,238,238}  % neutral panel
\definecolor{neutbord}{RGB}{170,170,170}
\definecolor{inkc}    {RGB}{ 30, 30, 30}  % strong ink (titles)
\definecolor{mutec}   {RGB}{120,120,120}  % muted text
\definecolor{focusc}  {RGB}{210, 70, 70}  % focus / highlight accent
\definecolor{gridc}   {RGB}{220,220,220}  % background grid

\tikzset{
  blog/.style={
    font=\sffamily\small\linespread{1.18}\selectfont,
    every node/.style={font=\sffamily\small\linespread{1.18}\selectfont}
  },
  % ─── existing node roles (unchanged) ───
  cls/.style={
    rectangle, rounded corners=2pt,
    draw=nodebord, line width=0.4pt,
    minimum height=8mm, minimum width=28mm,
    inner sep=3pt, fill=nodefill, align=center
  },
  abs/.style={cls, draw=absbord, fill=absfill, font=\sffamily\itshape\small},
  con/.style={cls, draw=conbord, fill=confill},
  imp/.style={cls, draw=impbord, fill=impfill},
  cli/.style={cls, draw=clibord, fill=clifill},
  hi/.style={cls, draw=hibord, fill=hifill, font=\sffamily\bfseries\small},
  % ─── new node roles ───
  prp/.style={cls, draw=prpbord, fill=prpfill},               % protocol / interface / state
  tealn/.style={cls, draw=tealbord, fill=tealfill},           % service / utility
  warn/.style={cls, draw=warnbord, fill=warnfill, font=\sffamily\bfseries\small},  % error state
  ok/.style={cls, draw=okbord, fill=okfill},                  % success state
  neut/.style={cls, draw=neutbord, fill=neutfill},            % neutral
  % ─── size variants ───
  clsSm/.style={cls, minimum height=6mm, minimum width=20mm, font=\sffamily\footnotesize, inner sep=2pt},
  clsLg/.style={cls, minimum height=11mm, minimum width=36mm, font=\sffamily\normalsize},
  absSm/.style={clsSm, draw=absbord, fill=absfill, font=\sffamily\itshape\footnotesize},
  conSm/.style={clsSm, draw=conbord, fill=confill},
  impSm/.style={clsSm, draw=impbord, fill=impfill},
  % ─── circle (existing) ───
  vx/.style={
    circle, draw=absbord, line width=0.4pt,
    minimum size=8mm, fill=absfill, inner sep=0pt
  },
  vxc/.style={vx, draw=conbord, fill=confill},
  vxh/.style={vx, draw=hibord, fill=hifill, font=\sffamily\bfseries},
  % new circle variants
  vxi/.style={vx, draw=impbord, fill=impfill},                % impl green
  vxp/.style={vx, draw=prpbord, fill=prpfill},                % protocol purple
  vxw/.style={vx, draw=warnbord, fill=warnfill},              % warn red
  vxn/.style={vx, draw=nodebord, fill=nodefill},              % neutral
  vxSm/.style={vx, minimum size=6mm, font=\sffamily\footnotesize},
  vxLg/.style={vx, minimum size=11mm, font=\sffamily\normalsize},
  % ─── decision diamond (existing) ───
  q/.style={
    diamond, draw=clibord, line width=0.4pt, fill=clifill,
    inner sep=2pt, aspect=2.4, align=center
  },
  % rounded callout / note (new)
  note/.style={
    rectangle, rounded corners=3pt,
    draw=clibord, line width=0.4pt, fill=clifill,
    inner sep=3pt, font=\sffamily\footnotesize, align=left
  },
  badge/.style={                                              % small inline marker
    rectangle, rounded corners=2pt,
    draw=hibord, line width=0.3pt, fill=hifill,
    inner sep=1.5pt, font=\sffamily\scriptsize\bfseries
  },
  % ─── edges (existing) ───
  inh/.style={                                  % UML inheritance
    -{Triangle[length=3mm,width=3mm,open]},
    draw=edgec, line width=0.5pt
  },
  agg/.style={                                  % UML aggregation
    {Diamond[length=2.5mm,width=2.5mm,open]}-{Stealth[length=2.2mm]},
    draw=edgec, line width=0.45pt
  },
  uses/.style={                                 % directed reference
    -{Stealth[length=2.2mm]}, draw=edgec, line width=0.5pt
  },
  arr/.style={                                  % plain arrow (graphs/trees)
    -{Stealth[length=2mm]}, draw=edgec, line width=0.45pt
  },
  edge/.style={                                 % undirected (trees)
    draw=edgec, line width=0.45pt
  },
  alt/.style={                                  % dashed: similar/alt
    -{Stealth[length=2mm]}, dashed, draw=edgesc, line width=0.45pt
  },
  ret/.style={                                  % dashed creates/returns
    -{Stealth[length=2mm]}, dashed, draw=edgesc, line width=0.45pt
  },
  thr/.style={                                  % threading link (red dashed)
    ->, dashed, draw=conbord, line width=0.5pt
  },
  % ─── new edges ───
  comp/.style={                                 % UML composition (filled diamond)
    {Diamond[length=2.5mm,width=2.5mm]}-{Stealth[length=2.2mm]},
    draw=edgec, line width=0.45pt
  },
  realize/.style={                              % UML realize (impl interface)
    -{Triangle[length=3mm,width=3mm,open]}, dashed,
    draw=edgec, line width=0.5pt
  },
  msgsync/.style={                              % synchronous message (sequence)
    -{Stealth[length=2.6mm]}, draw=edgec, line width=0.55pt
  },
  msgasync/.style={                             % asynchronous message
    -{Stealth[length=2.4mm,inset=0.4mm,fill=edgec]},
    draw=edgec, line width=0.5pt
  },
  msgret/.style={                               % return message (dashed)
    -{Stealth[length=2.2mm]}, dashed, draw=edgesc, line width=0.45pt
  },
  flow/.style={                                 % control flow (thicker, hi-vis)
    -{Stealth[length=2.4mm]}, draw=inkc, line width=0.7pt
  },
  focus/.style={                                % focus / hot path
    -{Stealth[length=2.4mm]}, draw=focusc, line width=0.8pt
  },
  selfloop/.style={                             % self loop on a node
    ->, draw=edgec, line width=0.45pt, in=120, out=60, loop, looseness=8
  },
  % ─── group / region (existing + new) ───
  group/.style={
    rectangle, rounded corners=4pt,
    draw=groupbord, line width=0.4pt,
    inner sep=10pt, fill=black!2
  },
  glabel/.style={font=\sffamily\bfseries\small, anchor=north west, inner sep=2pt},
  panel/.style={                                % stronger panel for layered diagrams
    rectangle, rounded corners=4pt,
    draw=neutbord, line width=0.45pt,
    inner sep=10pt, fill=neutfill!50
  },
  panelTitle/.style={font=\sffamily\bfseries\small, text=inkc, anchor=north west, inner sep=2pt},
  panelSub/.style={font=\sffamily\footnotesize, text=mutec, anchor=north west, inner sep=2pt},
  % ─── typography ───
  caption/.style={font=\sffamily\footnotesize\itshape, text=mutec, align=center},
  legendkey/.style={rectangle, minimum width=10mm, minimum height=3mm, inner sep=0pt, line width=0.3pt},
  % ─── grid / coordinate hints ───
  bggrid/.style={draw=gridc, line width=0.3pt, dotted},
}

% Korean font convenience macro — included by every Hangul diagram
\providecommand{\KOR}{}
           % core palette
%   % _design-sequence.tex — TikZ sequence diagram styles
%
% Mermaid는 *간단한* sequence diagram에 좋지만, 복잡한 경우 (여러 actor,
% 중첩 activation, alt/opt/loop frame, 자세한 라벨)는 auto-layout이 들쭉날쭉.
% 정밀 좌표가 필요할 때 TikZ로 그리는 표준.
%
% Companion to _design.tex — same palette and ink colors.
%
% Usage:
%   \documentclass[border=10pt,tikz]{standalone}
%   \usepackage{tikz}
%   \usetikzlibrary{positioning,arrows.meta,shapes,fit,calc,decorations.pathreplacing}
%   \begin{document}
%   % _design.tex — shared TikZ design system for all blog diagrams
% Inputs nothing on its own — meant to be \input{} from each diagram .tex
% AFTER \begin{document}, BEFORE \begin{tikzpicture}.
%
% Companion files:
%   _design-math.tex — math diagrams (PGFPlots, book-notes palette)
%
% Provides a unified visual language across GoF / DSA / EMC++ / Beautiful / 시리즈:
%   - Colors: muted, accessible, color-coded by role
%   - Node styles: consistent sizing, rounded, subtle borders
%   - Edge styles: UML-style inheritance, aggregation, plus dashed alternates
%   - Annotation: callouts, notes, badges
%
% Usage:
%   \documentclass[border=10pt,tikz]{standalone}
%   \usepackage{tikz}
%   \usetikzlibrary{positioning,arrows.meta,shapes,fit,backgrounds,calc,decorations.pathreplacing}
%   \begin{document}
%   \input{../_design.tex}     % adjust relative path
%   \begin{tikzpicture}[blog]  % apply 'blog' style
%   ...
%
% Backward-compat note: This file only ADDS to the public API. Existing styles
% (abs/con/imp/cli/hi/vx/q/inh/agg/uses/arr/edge/alt/ret/thr/group/glabel) are
% preserved. New helpers are documented below.

% ─── Core palette (fills paired with darker borders) ─────────
\definecolor{absfill}{RGB}{220,232,247}   % cool blue
\definecolor{absbord}{RGB}{72,118,182}
\definecolor{confill}{RGB}{249,222,222}   % warm red
\definecolor{conbord}{RGB}{197,93,93}
\definecolor{impfill}{RGB}{223,238,222}   % soft green
\definecolor{impbord}{RGB}{99,154,93}
\definecolor{clifill}{RGB}{251,239,205}   % pale yellow
\definecolor{clibord}{RGB}{200,167,72}
\definecolor{hifill} {RGB}{253,224,200}   % accent orange
\definecolor{hibord} {RGB}{210,135,68}
\definecolor{nodefill}{RGB}{250,250,250}
\definecolor{nodebord}{RGB}{120,120,120}
\definecolor{edgec}  {RGB}{80,80,80}
\definecolor{edgesc} {RGB}{145,145,145}
\definecolor{groupbord}{RGB}{200,200,200}

% ─── Extended palette (new) ──────────────────────────────────
\definecolor{prpfill}{RGB}{233,224,247}   % muted purple (state / interface)
\definecolor{prpbord}{RGB}{132,103,182}
\definecolor{tealfill}{RGB}{215,235,235}  % teal (services / utility)
\definecolor{tealbord}{RGB}{ 78,135,140}
\definecolor{warnfill}{RGB}{251,229,229}  % deeper warn (error states)
\definecolor{warnbord}{RGB}{180, 55, 55}
\definecolor{okfill}  {RGB}{220,240,225}  % success
\definecolor{okbord}  {RGB}{ 75,140, 90}
\definecolor{neutfill}{RGB}{238,238,238}  % neutral panel
\definecolor{neutbord}{RGB}{170,170,170}
\definecolor{inkc}    {RGB}{ 30, 30, 30}  % strong ink (titles)
\definecolor{mutec}   {RGB}{120,120,120}  % muted text
\definecolor{focusc}  {RGB}{210, 70, 70}  % focus / highlight accent
\definecolor{gridc}   {RGB}{220,220,220}  % background grid

\tikzset{
  blog/.style={
    font=\sffamily\small\linespread{1.18}\selectfont,
    every node/.style={font=\sffamily\small\linespread{1.18}\selectfont}
  },
  % ─── existing node roles (unchanged) ───
  cls/.style={
    rectangle, rounded corners=2pt,
    draw=nodebord, line width=0.4pt,
    minimum height=8mm, minimum width=28mm,
    inner sep=3pt, fill=nodefill, align=center
  },
  abs/.style={cls, draw=absbord, fill=absfill, font=\sffamily\itshape\small},
  con/.style={cls, draw=conbord, fill=confill},
  imp/.style={cls, draw=impbord, fill=impfill},
  cli/.style={cls, draw=clibord, fill=clifill},
  hi/.style={cls, draw=hibord, fill=hifill, font=\sffamily\bfseries\small},
  % ─── new node roles ───
  prp/.style={cls, draw=prpbord, fill=prpfill},               % protocol / interface / state
  tealn/.style={cls, draw=tealbord, fill=tealfill},           % service / utility
  warn/.style={cls, draw=warnbord, fill=warnfill, font=\sffamily\bfseries\small},  % error state
  ok/.style={cls, draw=okbord, fill=okfill},                  % success state
  neut/.style={cls, draw=neutbord, fill=neutfill},            % neutral
  % ─── size variants ───
  clsSm/.style={cls, minimum height=6mm, minimum width=20mm, font=\sffamily\footnotesize, inner sep=2pt},
  clsLg/.style={cls, minimum height=11mm, minimum width=36mm, font=\sffamily\normalsize},
  absSm/.style={clsSm, draw=absbord, fill=absfill, font=\sffamily\itshape\footnotesize},
  conSm/.style={clsSm, draw=conbord, fill=confill},
  impSm/.style={clsSm, draw=impbord, fill=impfill},
  % ─── circle (existing) ───
  vx/.style={
    circle, draw=absbord, line width=0.4pt,
    minimum size=8mm, fill=absfill, inner sep=0pt
  },
  vxc/.style={vx, draw=conbord, fill=confill},
  vxh/.style={vx, draw=hibord, fill=hifill, font=\sffamily\bfseries},
  % new circle variants
  vxi/.style={vx, draw=impbord, fill=impfill},                % impl green
  vxp/.style={vx, draw=prpbord, fill=prpfill},                % protocol purple
  vxw/.style={vx, draw=warnbord, fill=warnfill},              % warn red
  vxn/.style={vx, draw=nodebord, fill=nodefill},              % neutral
  vxSm/.style={vx, minimum size=6mm, font=\sffamily\footnotesize},
  vxLg/.style={vx, minimum size=11mm, font=\sffamily\normalsize},
  % ─── decision diamond (existing) ───
  q/.style={
    diamond, draw=clibord, line width=0.4pt, fill=clifill,
    inner sep=2pt, aspect=2.4, align=center
  },
  % rounded callout / note (new)
  note/.style={
    rectangle, rounded corners=3pt,
    draw=clibord, line width=0.4pt, fill=clifill,
    inner sep=3pt, font=\sffamily\footnotesize, align=left
  },
  badge/.style={                                              % small inline marker
    rectangle, rounded corners=2pt,
    draw=hibord, line width=0.3pt, fill=hifill,
    inner sep=1.5pt, font=\sffamily\scriptsize\bfseries
  },
  % ─── edges (existing) ───
  inh/.style={                                  % UML inheritance
    -{Triangle[length=3mm,width=3mm,open]},
    draw=edgec, line width=0.5pt
  },
  agg/.style={                                  % UML aggregation
    {Diamond[length=2.5mm,width=2.5mm,open]}-{Stealth[length=2.2mm]},
    draw=edgec, line width=0.45pt
  },
  uses/.style={                                 % directed reference
    -{Stealth[length=2.2mm]}, draw=edgec, line width=0.5pt
  },
  arr/.style={                                  % plain arrow (graphs/trees)
    -{Stealth[length=2mm]}, draw=edgec, line width=0.45pt
  },
  edge/.style={                                 % undirected (trees)
    draw=edgec, line width=0.45pt
  },
  alt/.style={                                  % dashed: similar/alt
    -{Stealth[length=2mm]}, dashed, draw=edgesc, line width=0.45pt
  },
  ret/.style={                                  % dashed creates/returns
    -{Stealth[length=2mm]}, dashed, draw=edgesc, line width=0.45pt
  },
  thr/.style={                                  % threading link (red dashed)
    ->, dashed, draw=conbord, line width=0.5pt
  },
  % ─── new edges ───
  comp/.style={                                 % UML composition (filled diamond)
    {Diamond[length=2.5mm,width=2.5mm]}-{Stealth[length=2.2mm]},
    draw=edgec, line width=0.45pt
  },
  realize/.style={                              % UML realize (impl interface)
    -{Triangle[length=3mm,width=3mm,open]}, dashed,
    draw=edgec, line width=0.5pt
  },
  msgsync/.style={                              % synchronous message (sequence)
    -{Stealth[length=2.6mm]}, draw=edgec, line width=0.55pt
  },
  msgasync/.style={                             % asynchronous message
    -{Stealth[length=2.4mm,inset=0.4mm,fill=edgec]},
    draw=edgec, line width=0.5pt
  },
  msgret/.style={                               % return message (dashed)
    -{Stealth[length=2.2mm]}, dashed, draw=edgesc, line width=0.45pt
  },
  flow/.style={                                 % control flow (thicker, hi-vis)
    -{Stealth[length=2.4mm]}, draw=inkc, line width=0.7pt
  },
  focus/.style={                                % focus / hot path
    -{Stealth[length=2.4mm]}, draw=focusc, line width=0.8pt
  },
  selfloop/.style={                             % self loop on a node
    ->, draw=edgec, line width=0.45pt, in=120, out=60, loop, looseness=8
  },
  % ─── group / region (existing + new) ───
  group/.style={
    rectangle, rounded corners=4pt,
    draw=groupbord, line width=0.4pt,
    inner sep=10pt, fill=black!2
  },
  glabel/.style={font=\sffamily\bfseries\small, anchor=north west, inner sep=2pt},
  panel/.style={                                % stronger panel for layered diagrams
    rectangle, rounded corners=4pt,
    draw=neutbord, line width=0.45pt,
    inner sep=10pt, fill=neutfill!50
  },
  panelTitle/.style={font=\sffamily\bfseries\small, text=inkc, anchor=north west, inner sep=2pt},
  panelSub/.style={font=\sffamily\footnotesize, text=mutec, anchor=north west, inner sep=2pt},
  % ─── typography ───
  caption/.style={font=\sffamily\footnotesize\itshape, text=mutec, align=center},
  legendkey/.style={rectangle, minimum width=10mm, minimum height=3mm, inner sep=0pt, line width=0.3pt},
  % ─── grid / coordinate hints ───
  bggrid/.style={draw=gridc, line width=0.3pt, dotted},
}

% Korean font convenience macro — included by every Hangul diagram
\providecommand{\KOR}{}
           % core palette
%   % _design-sequence.tex — TikZ sequence diagram styles
%
% Mermaid는 *간단한* sequence diagram에 좋지만, 복잡한 경우 (여러 actor,
% 중첩 activation, alt/opt/loop frame, 자세한 라벨)는 auto-layout이 들쭉날쭉.
% 정밀 좌표가 필요할 때 TikZ로 그리는 표준.
%
% Companion to _design.tex — same palette and ink colors.
%
% Usage:
%   \documentclass[border=10pt,tikz]{standalone}
%   \usepackage{tikz}
%   \usetikzlibrary{positioning,arrows.meta,shapes,fit,calc,decorations.pathreplacing}
%   \begin{document}
%   \input{../_design.tex}           % core palette
%   \input{../_design-sequence.tex}  % sequence-specific
%   \begin{tikzpicture}[blog, seq]
%     % Actors (top)
%     \node[actor]   (a) at (0, 0)  {Client};
%     \node[actor]   (b) at (4, 0)  {Server};
%     \node[actor]   (c) at (8, 0)  {DB};
%
%     % Lifelines (dashed vertical, automatic length via target node)
%     \draw[lifeline] (a) -- ++(0,-7);
%     \draw[lifeline] (b) -- ++(0,-7);
%     \draw[lifeline] (c) -- ++(0,-7);
%
%     % Activation boxes (on lifeline)
%     \node[activation, minimum height=14mm] at (b |- 0,-2.2) {};
%
%     % Messages (top-to-bottom time axis)
%     \draw[msgsync] (a |- 0,-1.5) -- node[msglabel]{request} (b |- 0,-1.5);
%     \draw[msgsync] (b |- 0,-2.5) -- node[msglabel]{query}   (c |- 0,-2.5);
%     \draw[msgret]  (c |- 0,-3.2) -- node[msglabel]{rows}    (b |- 0,-3.2);
%     \draw[msgret]  (b |- 0,-4.0) -- node[msglabel]{response}(a |- 0,-4.0);
%
%     % Frame (alt / opt / loop)
%     \node[frame, fit={(2,-5) (6,-6.5)}, label={[framelabel]north west:alt}] {};
%   \end{tikzpicture}

\tikzset{
  % Apply 'seq' style on top of 'blog' to get sequence-friendly defaults.
  seq/.style={
    every node/.append style={font=\sffamily\small},
  },
  % ─── Actors (top of diagram) ────────────────────────────
  actor/.style={
    rectangle, rounded corners=2pt,
    draw=nodebord, line width=0.45pt,
    minimum height=8mm, minimum width=22mm,
    inner sep=3pt, fill=nodefill, align=center,
    font=\sffamily\bfseries\small,
  },
  actorAlt/.style={actor, draw=absbord, fill=absfill},     % alternate accent
  actorWarn/.style={actor, draw=warnbord, fill=warnfill},
  % ─── Lifeline (dashed vertical) ─────────────────────────
  lifeline/.style={
    dashed, draw=mutec, line width=0.4pt,
  },
  % ─── Activation box (process running) ───────────────────
  activation/.style={
    rectangle, fill=absfill, draw=absbord, line width=0.35pt,
    minimum width=2.4mm, anchor=north,
    inner sep=0pt,
  },
  activationAlt/.style={activation, fill=hifill, draw=hibord},  % nested
  % ─── Messages ────────────────────────────────────────────
  % msgsync / msgasync / msgret already defined in _design.tex; we add labels & specials.
  msglabel/.style={
    font=\sffamily\footnotesize, text=inkc,
    midway, above, fill=white, fill opacity=0.85, text opacity=1,
    inner sep=2pt,
  },
  msglabelBelow/.style={msglabel, above=false, below},
  msgself/.style={                                          % self-message (loop on lifeline)
    -{Stealth[length=2.2mm]}, draw=edgec, line width=0.5pt,
  },
  msgcreate/.style={                                        % object creation (sets up new lifeline)
    -{Stealth[length=2.4mm]}, dashed, draw=edgec, line width=0.5pt,
  },
  msgdestroy/.style={                                       % destruction (cross at lifeline end)
    -{Stealth[length=2mm]}, draw=warnbord, line width=0.55pt,
  },
  % destruction cross marker
  destroyX/.pic={
    \draw[warnbord, line width=0.6pt] (-1.2mm,-1.2mm) -- (1.2mm,1.2mm);
    \draw[warnbord, line width=0.6pt] (-1.2mm, 1.2mm) -- (1.2mm,-1.2mm);
  },
  % ─── Frame (alt / opt / loop / par) ─────────────────────
  frame/.style={
    rectangle,
    draw=mutec, line width=0.4pt, dashed,
    inner sep=4pt,
    rounded corners=2pt,
  },
  framelabel/.style={
    font=\sffamily\bfseries\footnotesize, text=inkc,
    fill=neutfill, draw=mutec, line width=0.3pt,
    inner sep=2pt,
    rounded corners=1pt,
  },
  frameAltSep/.style={                                     % horizontal separator inside alt frame
    dashed, draw=mutec, line width=0.4pt,
  },
  % ─── Note / comment attached to a message or lifeline ───
  seqnote/.style={
    rectangle, rounded corners=2pt,
    draw=clibord, line width=0.4pt, fill=clifill,
    inner sep=3pt, font=\sffamily\footnotesize, align=left,
  },
  seqnoteAlt/.style={seqnote, draw=tealbord, fill=tealfill},
  % ─── Time-axis hints (rarely needed; use sparingly) ─────
  timearrow/.style={
    -{Stealth[length=2.4mm]}, draw=mutec, line width=0.3pt,
  },
  timetick/.style={
    draw=mutec, line width=0.3pt,
  },
}
  % sequence-specific
%   \begin{tikzpicture}[blog, seq]
%     % Actors (top)
%     \node[actor]   (a) at (0, 0)  {Client};
%     \node[actor]   (b) at (4, 0)  {Server};
%     \node[actor]   (c) at (8, 0)  {DB};
%
%     % Lifelines (dashed vertical, automatic length via target node)
%     \draw[lifeline] (a) -- ++(0,-7);
%     \draw[lifeline] (b) -- ++(0,-7);
%     \draw[lifeline] (c) -- ++(0,-7);
%
%     % Activation boxes (on lifeline)
%     \node[activation, minimum height=14mm] at (b |- 0,-2.2) {};
%
%     % Messages (top-to-bottom time axis)
%     \draw[msgsync] (a |- 0,-1.5) -- node[msglabel]{request} (b |- 0,-1.5);
%     \draw[msgsync] (b |- 0,-2.5) -- node[msglabel]{query}   (c |- 0,-2.5);
%     \draw[msgret]  (c |- 0,-3.2) -- node[msglabel]{rows}    (b |- 0,-3.2);
%     \draw[msgret]  (b |- 0,-4.0) -- node[msglabel]{response}(a |- 0,-4.0);
%
%     % Frame (alt / opt / loop)
%     \node[frame, fit={(2,-5) (6,-6.5)}, label={[framelabel]north west:alt}] {};
%   \end{tikzpicture}

\tikzset{
  % Apply 'seq' style on top of 'blog' to get sequence-friendly defaults.
  seq/.style={
    every node/.append style={font=\sffamily\small},
  },
  % ─── Actors (top of diagram) ────────────────────────────
  actor/.style={
    rectangle, rounded corners=2pt,
    draw=nodebord, line width=0.45pt,
    minimum height=8mm, minimum width=22mm,
    inner sep=3pt, fill=nodefill, align=center,
    font=\sffamily\bfseries\small,
  },
  actorAlt/.style={actor, draw=absbord, fill=absfill},     % alternate accent
  actorWarn/.style={actor, draw=warnbord, fill=warnfill},
  % ─── Lifeline (dashed vertical) ─────────────────────────
  lifeline/.style={
    dashed, draw=mutec, line width=0.4pt,
  },
  % ─── Activation box (process running) ───────────────────
  activation/.style={
    rectangle, fill=absfill, draw=absbord, line width=0.35pt,
    minimum width=2.4mm, anchor=north,
    inner sep=0pt,
  },
  activationAlt/.style={activation, fill=hifill, draw=hibord},  % nested
  % ─── Messages ────────────────────────────────────────────
  % msgsync / msgasync / msgret already defined in _design.tex; we add labels & specials.
  msglabel/.style={
    font=\sffamily\footnotesize, text=inkc,
    midway, above, fill=white, fill opacity=0.85, text opacity=1,
    inner sep=2pt,
  },
  msglabelBelow/.style={msglabel, above=false, below},
  msgself/.style={                                          % self-message (loop on lifeline)
    -{Stealth[length=2.2mm]}, draw=edgec, line width=0.5pt,
  },
  msgcreate/.style={                                        % object creation (sets up new lifeline)
    -{Stealth[length=2.4mm]}, dashed, draw=edgec, line width=0.5pt,
  },
  msgdestroy/.style={                                       % destruction (cross at lifeline end)
    -{Stealth[length=2mm]}, draw=warnbord, line width=0.55pt,
  },
  % destruction cross marker
  destroyX/.pic={
    \draw[warnbord, line width=0.6pt] (-1.2mm,-1.2mm) -- (1.2mm,1.2mm);
    \draw[warnbord, line width=0.6pt] (-1.2mm, 1.2mm) -- (1.2mm,-1.2mm);
  },
  % ─── Frame (alt / opt / loop / par) ─────────────────────
  frame/.style={
    rectangle,
    draw=mutec, line width=0.4pt, dashed,
    inner sep=4pt,
    rounded corners=2pt,
  },
  framelabel/.style={
    font=\sffamily\bfseries\footnotesize, text=inkc,
    fill=neutfill, draw=mutec, line width=0.3pt,
    inner sep=2pt,
    rounded corners=1pt,
  },
  frameAltSep/.style={                                     % horizontal separator inside alt frame
    dashed, draw=mutec, line width=0.4pt,
  },
  % ─── Note / comment attached to a message or lifeline ───
  seqnote/.style={
    rectangle, rounded corners=2pt,
    draw=clibord, line width=0.4pt, fill=clifill,
    inner sep=3pt, font=\sffamily\footnotesize, align=left,
  },
  seqnoteAlt/.style={seqnote, draw=tealbord, fill=tealfill},
  % ─── Time-axis hints (rarely needed; use sparingly) ─────
  timearrow/.style={
    -{Stealth[length=2.4mm]}, draw=mutec, line width=0.3pt,
  },
  timetick/.style={
    draw=mutec, line width=0.3pt,
  },
}
  % sequence-specific
%   \begin{tikzpicture}[blog, seq]
%     % Actors (top)
%     \node[actor]   (a) at (0, 0)  {Client};
%     \node[actor]   (b) at (4, 0)  {Server};
%     \node[actor]   (c) at (8, 0)  {DB};
%
%     % Lifelines (dashed vertical, automatic length via target node)
%     \draw[lifeline] (a) -- ++(0,-7);
%     \draw[lifeline] (b) -- ++(0,-7);
%     \draw[lifeline] (c) -- ++(0,-7);
%
%     % Activation boxes (on lifeline)
%     \node[activation, minimum height=14mm] at (b |- 0,-2.2) {};
%
%     % Messages (top-to-bottom time axis)
%     \draw[msgsync] (a |- 0,-1.5) -- node[msglabel]{request} (b |- 0,-1.5);
%     \draw[msgsync] (b |- 0,-2.5) -- node[msglabel]{query}   (c |- 0,-2.5);
%     \draw[msgret]  (c |- 0,-3.2) -- node[msglabel]{rows}    (b |- 0,-3.2);
%     \draw[msgret]  (b |- 0,-4.0) -- node[msglabel]{response}(a |- 0,-4.0);
%
%     % Frame (alt / opt / loop)
%     \node[frame, fit={(2,-5) (6,-6.5)}, label={[framelabel]north west:alt}] {};
%   \end{tikzpicture}

\tikzset{
  % Apply 'seq' style on top of 'blog' to get sequence-friendly defaults.
  seq/.style={
    every node/.append style={font=\sffamily\small},
  },
  % ─── Actors (top of diagram) ────────────────────────────
  actor/.style={
    rectangle, rounded corners=2pt,
    draw=nodebord, line width=0.45pt,
    minimum height=8mm, minimum width=22mm,
    inner sep=3pt, fill=nodefill, align=center,
    font=\sffamily\bfseries\small,
  },
  actorAlt/.style={actor, draw=absbord, fill=absfill},     % alternate accent
  actorWarn/.style={actor, draw=warnbord, fill=warnfill},
  % ─── Lifeline (dashed vertical) ─────────────────────────
  lifeline/.style={
    dashed, draw=mutec, line width=0.4pt,
  },
  % ─── Activation box (process running) ───────────────────
  activation/.style={
    rectangle, fill=absfill, draw=absbord, line width=0.35pt,
    minimum width=2.4mm, anchor=north,
    inner sep=0pt,
  },
  activationAlt/.style={activation, fill=hifill, draw=hibord},  % nested
  % ─── Messages ────────────────────────────────────────────
  % msgsync / msgasync / msgret already defined in _design.tex; we add labels & specials.
  msglabel/.style={
    font=\sffamily\footnotesize, text=inkc,
    midway, above, fill=white, fill opacity=0.85, text opacity=1,
    inner sep=2pt,
  },
  msglabelBelow/.style={msglabel, above=false, below},
  msgself/.style={                                          % self-message (loop on lifeline)
    -{Stealth[length=2.2mm]}, draw=edgec, line width=0.5pt,
  },
  msgcreate/.style={                                        % object creation (sets up new lifeline)
    -{Stealth[length=2.4mm]}, dashed, draw=edgec, line width=0.5pt,
  },
  msgdestroy/.style={                                       % destruction (cross at lifeline end)
    -{Stealth[length=2mm]}, draw=warnbord, line width=0.55pt,
  },
  % destruction cross marker
  destroyX/.pic={
    \draw[warnbord, line width=0.6pt] (-1.2mm,-1.2mm) -- (1.2mm,1.2mm);
    \draw[warnbord, line width=0.6pt] (-1.2mm, 1.2mm) -- (1.2mm,-1.2mm);
  },
  % ─── Frame (alt / opt / loop / par) ─────────────────────
  frame/.style={
    rectangle,
    draw=mutec, line width=0.4pt, dashed,
    inner sep=4pt,
    rounded corners=2pt,
  },
  framelabel/.style={
    font=\sffamily\bfseries\footnotesize, text=inkc,
    fill=neutfill, draw=mutec, line width=0.3pt,
    inner sep=2pt,
    rounded corners=1pt,
  },
  frameAltSep/.style={                                     % horizontal separator inside alt frame
    dashed, draw=mutec, line width=0.4pt,
  },
  % ─── Note / comment attached to a message or lifeline ───
  seqnote/.style={
    rectangle, rounded corners=2pt,
    draw=clibord, line width=0.4pt, fill=clifill,
    inner sep=3pt, font=\sffamily\footnotesize, align=left,
  },
  seqnoteAlt/.style={seqnote, draw=tealbord, fill=tealfill},
  % ─── Time-axis hints (rarely needed; use sparingly) ─────
  timearrow/.style={
    -{Stealth[length=2.4mm]}, draw=mutec, line width=0.3pt,
  },
  timetick/.style={
    draw=mutec, line width=0.3pt,
  },
}

\begin{tikzpicture}[blog, seq]

% Actors (top)
\node[actor]    (p)  at (0,   0) {Primary CPU\\(EL1)};
\node[actorAlt] (b)  at (6,   0) {BL31\\(EL3)};
\node[actorWarn] (sc) at (12,  0) {SoC reset\\controller};
\node[actor]    (s)  at (17,  0) {Secondary\\CPU};

% Lifelines
\draw[lifeline] (p)  -- ++(0,-9.2);
\draw[lifeline] (b)  -- ++(0,-9.2);
\draw[lifeline] (sc) -- ++(0,-9.2);
\draw[lifeline] (s)  -- ++(0,-9.2);

% Messages — keep y separated by >=0.9 so labels do not collide

% 1. SMC trap
\draw[msgsync] (p |- 0,-1.2) -- node[msglabel]{\texttt{smc \#0} : CPU\_ON(target, entry, ctx)} (b |- 0,-1.2);

% 2. validate inputs (self-message on BL31)
\draw[msgsync] (b |- 0,-2.3) -- ++(1.2,0) |- node[msglabel, pos=0.25]{validate MPIDR / entry} (b |- 0,-3.0);

% 3. pwr_domain_on → SoC SRC
\draw[msgsync] (b |- 0,-4.0) -- node[msglabel]{pwr\_domain\_on()\\write SRC reset register} (sc |- 0,-4.0);

% 4. reset deassert → secondary CPU starts
\draw[msgasync] (sc |- 0,-5.0) -- node[msglabel]{reset deassert} (s |- 0,-5.0);

% 5. SMC return (PSCI_SUCCESS)
\draw[msgret] (b |- 0,-6.0) -- node[msglabel]{PSCI\_SUCCESS} (p |- 0,-6.0);

% 6. Secondary enters BL31 warm boot
\draw[msgsync] (s |- 0,-7.0) -- ++(-1.5,0) |- node[msglabel, pos=0.25]{warm boot entry} (s |- 0,-7.7);

% 7. Secondary jumps to caller-provided entry
\draw[msgsync] (s |- 0,-8.7) -- ++(-1.5,0) |- node[msglabel, pos=0.25]{br x4 = entry\\x0 = ctx} (s |- 0,-9.0);

% Note explaining async nature (place to the left, away from message labels)
\node[seqnote, anchor=north west] at (-2.5, -7.0)
  {SMC는 \texttt{PSCI\_SUCCESS}로\\[2pt]즉시 return하지만\\[2pt]secondary는 비동기로\\[2pt]entry에 도착한다.};

\end{tikzpicture}
\end{document}
